\documentclass[11pt,a4paper]{article}

% Minimal, pdflatex-friendly setup
\usepackage[T1]{fontenc}
\usepackage[utf8]{inputenc}
\usepackage{lmodern}
\usepackage[a4paper,margin=1in]{geometry}
\usepackage{amsmath,amssymb}
\usepackage{hyperref}
\usepackage{verbatim} % for code blocks

\title{SageMath Quick Notes: Derivatives, Rendering to LaTeX, and Basic Plotting}
\author{}
\date{}

\begin{document}
	\maketitle
	To start it :
	
	\begin{verbatim}
		sage -n jupyter
	\end{verbatim}
	
	\section*{Conventions}
	We use capital $I$ for the imaginary unit in Sage, i.e.\ $I=\sqrt{-1}$. 
	If you prefer engineering $j$, you can alias \verb|j = I|.
	
	\section{Variables and Basic Expressions}
	\begin{verbatim}
		# Declare symbols
		t, x, s, w = var('t x s w')
		
		# Define a few expressions
		f  = exp(-2*t)*cos(3*t)              # e^{-2t} cos(3t)
		g  = 20*log(sqrt(1 + 1/w^2), 10)     # 20 log_10(sqrt(1 + 1/w^2))
		z  = exp((2 + 3*I)*t)                # e^{(2+3i)t}
	\end{verbatim}
	
	Displayed mathematics we are modeling with Sage:
	$$
	f(t)=e^{-2t}\cos(3t),\qquad
	G(\omega)=20\log_{10}\Big(\sqrt{1+\tfrac{1}{\omega^2}}\Big),\qquad
	z(t)=e^{(2+3i)t}.
	$$
	
	\section{Derivatives (warm-up)}
	\begin{verbatim}
		df_dt   = diff(f, t)                 # derivative of f w.r.t. t
		d2z_dt2 = diff(z, t, 2)              # second derivative of z
		dG_dw   = diff(g, w)                 # derivative of G with respect to w
	\end{verbatim}
	
	Example target form:
	$$
	\frac{d}{dt}\big(e^{-2t}\cos 3t\big)
	= e^{-2t}\big(-2\cos(3t) - 3\sin(3t)\big).
	$$
	
	\section{Logs and Bases}
	Natural log is $\ln x$ via \verb|log(x)|. 
	Base-10 is $\log_{10}x$ via \verb|log(x,10)|.
	\begin{verbatim}
		lnx   = log(x)        # ln(x)
		lgx   = log(x, 10)    # log_10(x)
		lgb   = log(x, 2)     # log_2(x)
	\end{verbatim}
	In maths:
	$$
	\log x \equiv \ln x,\qquad \log_{10}x = \frac{\ln x}{\ln 10}.
	$$
	
	\section{Complex Unit: I (and alias j)}
	Sage uses $I$. You can alias $j$ to look like EE notes:
	\begin{verbatim}
		j = I
		zj = exp((2 + 3*j)*t)  # identical to using I under the hood
	\end{verbatim}
	Mathematically:
	$$
	\frac{d}{dt}e^{(2+3i)t}=(2+3i)e^{(2+3i)t}.
	$$
	
	\section{Pretty Rendering vs. LaTeX Source}
	Use \verb|show(expr)| for nicely rendered maths inside Sage/Jupyter.  
	Use \verb|latex(expr)| to get the LaTeX code you can paste into your document.
	
	\subsection*{Examples}
	\begin{verbatim}
		# Pretty in-notebook display
		show(f)        # shows e^{-2 t} cos(3 t) typeset
		show(df_dt)    # shows derivative typeset
		show(z)
		
		# Get LaTeX source strings (copy-paste into your .tex)
		latex_f    = latex(f)       
		latex_dfdt = latex(df_dt)
		latex_z    = latex(z)
		
		print(latex_f)
		print(latex_dfdt)
		print(latex_z)
	\end{verbatim}
	
	Expected outputs:
	$$
	\texttt{latex(f)}:\ e^{-2 t}\cos(3 t),\qquad
	\texttt{latex(z)}:\ e^{(2 + 3 i) t}.
	$$
	
	\section{Basic Plotting (2D)}
	\subsection*{Single function}
	\begin{verbatim}
		p1 = plot(f, (t, -2, 2), axes_labels=('t','f(t)'), legend_label='f(t)')
		p1.show()
	\end{verbatim}
	
	\subsection*{Overlay multiple curves}
	\begin{verbatim}
		df_dt = diff(f, t)
		p2 = plot(f, (t, -2, 2), legend_label='f(t)')
		p3 = plot(df_dt, (t, -2, 2), color='red', legend_label="f'(t)")
		(p2 + p3).show()
	\end{verbatim}
	
	\subsection*{Bode-like magnitude}
	For $G(\omega)=20\log_{10}(\sqrt{1+1/\omega^2})$:
	\begin{verbatim}
		pG = plot(20*log(sqrt(1 + 1/w^2), 10), (w, 0.1, 10),
		axes_labels=('w','G(w) [dB]'), legend_label='G(w)')
		pG.show()
	\end{verbatim}
	
	\section{Real and Imag Parts of Complex Signals}
	Let $z(t)=e^{(2+3i)t}=e^{2t}(\cos 3t + i\sin 3t)$.  
	\begin{verbatim}
		Re = real_part(z)   # e^(2*t)*cos(3*t)
		Im = imag_part(z)   # e^(2*t)*sin(3*t)
		
		pR = plot(Re, (t, 0, 2), legend_label='Re z(t)')
		pI = plot(Im, (t, 0, 2), color='green', legend_label='Im z(t)')
		(pR + pI).show()
	\end{verbatim}
	In maths:
	$$
	\Re z(t)=e^{2t}\cos(3t),\qquad \Im z(t)=e^{2t}\sin(3t).
	$$
	
	\section{End-to-End Workflow}
	Goal: differentiate $f(t)=e^{-2t}\cos(3t)$, render nicely, and plot with its derivative.
	\begin{verbatim}
		# 1) Define
		t = var('t')
		f = exp(-2*t)*cos(3*t)
		
		# 2) Differentiate
		df = diff(f, t)
		
		# 3) Pretty-display and LaTeX source
		show(f); show(df)
		print(latex(f))
		print(latex(df))
		
		# 4) Plot
		pf  = plot(f,  (t, -2, 2), legend_label='f(t)')
		pdf = plot(df, (t, -2, 2), color='red', legend_label="f'(t)")
		(pf + pdf).show()
	\end{verbatim}
	
	Expected derivative:
	$$
	f'(t)=e^{-2t}\big(-2\cos(3t)-3\sin(3t)\big).
	$$
	
	\bigskip
	\noindent\textbf{Notes.}
	\begin{itemize}
		\item Use \verb|simplify_trig()| or \verb|full_simplify()| when expressions look messy.
		\item Prefer \verb|show(expr)| for quick rendering; use \verb|latex(expr)| when you need to paste into documents.
		\item For engineering style, alias \verb|j = I|.
	\end{itemize}
	
\end{document}
\right) 